\section{Limitations and Conclusion}
% provenance: OUTLINE.md status column; jrn_01KXT4NN8Z8EG0SZYGB65H2T5J (venue); chk T1.1 open item.

\paragraph{Limitations.}
The boundaries of the evaluation are stated next to the results they bound, at the end of
Section~\ref{sec:prelim}: one feeder and a single solver, one PV penetration level at the
conformant scope envelope, modeled rather than measured retained authority, workstation-only
overhead, in-process availability injection, untested interoperability, an unmeasured
stealth-constrained attacker, and no access to the 2023 edition of the standard. Two further
boundaries belong to the model. The authorization graphs used by the engine are synthetic, because
real utility authorization graphs are not public; we mitigate this with a released parameterized
generator and by ablating access-model permissiveness rather than fixing it. And our claims about
revocation rest on the Sandia analysis rather than a clause-level reading of the edition in force.

\paragraph{Conclusion.}
Credential compromise in IEEE~2030.5 is not well described by how many nodes an adversary reaches.
It is described by how much physical capacity those nodes command, how long the adversary keeps
it, and what the feeder does in response. Our central empirical finding is about the \emph{shape} of the
authorization an operator grants, not about how tightly it is drawn. \textbf{On the
IEEE~8500-node feeder, a reactive authorization that permits zero absorption does not contain a
credential compromise, and one that mandates a floor on absorption does.} An adversary needing no
injection authority at all --- commanding zero, which sits inside any magnitude bound --- drives
the feeder past an overvoltage trip threshold in every seed, because withdrawing the volt-var
absorption is itself the attack. An adversary confined instead to a band holding at least about
three quarters of the conformant Category~B absorption trips no seed. Least privilege works here,
but only when the bound is a floor rather than a ceiling.

We expected the control primitive to be what separated these cases and it is not: we ran the sweep
under both a fixed-setpoint override and a bounded volt-var curve, and over the band an operator
would authorize the two are indistinguishable on the trip outcome. Reporting a primitive-level
separation would have described the shape of the two authorized sets we chose rather than anything
the feeder did, so we do not. Whether the standard's function sets can express an absorption floor
for a fixed setpoint is a question about IEEE~2030.5 that we could not answer without the edition
in force, and we leave it flagged as a conjecture in Section~\ref{sec:prelim}.

Two things follow. The concrete one is a recommendation to the standard's users: authorize a
reactive band that excludes zero absorption, by whichever function set carries it. The
second is that scope discipline remains insufficient on its own at the envelopes deployments
actually grant today, so bounding the \emph{duration} of retained authority still matters, and
that is what attested short-lived credentials with session and command-effect enforcement
provide. The price is stated rather than hidden: containment latency is the credential
lifetime divided by the per-cycle attestation detection probability, not the lifetime itself, and a
gateway-proxy tier is the only one of our three deployment tiers that resolves the secure-time and
constrained-hardware constraints the mechanism inherits from the problem it addresses. These are
pilot results on one feeder; the pre-registered confirmatory sweep is what would generalise them.
